% Options for packages loaded elsewhere
\PassOptionsToPackage{unicode}{hyperref}
\PassOptionsToPackage{hyphens}{url}
\PassOptionsToPackage{dvipsnames,svgnames,x11names}{xcolor}
\documentclass[
]{article}
\usepackage{xcolor}
\usepackage[margin = 0.5in]{geometry}
\usepackage{amsmath,amssymb}
\setcounter{secnumdepth}{-\maxdimen} % remove section numbering
\usepackage{iftex}
\ifPDFTeX
  \usepackage[T1]{fontenc}
  \usepackage[utf8]{inputenc}
  \usepackage{textcomp} % provide euro and other symbols
\else % if luatex or xetex
  \usepackage{unicode-math} % this also loads fontspec
  \defaultfontfeatures{Scale=MatchLowercase}
  \defaultfontfeatures[\rmfamily]{Ligatures=TeX,Scale=1}
\fi
\usepackage{lmodern}
\ifPDFTeX\else
  % xetex/luatex font selection
\fi
% Use upquote if available, for straight quotes in verbatim environments
\IfFileExists{upquote.sty}{\usepackage{upquote}}{}
\IfFileExists{microtype.sty}{% use microtype if available
  \usepackage[]{microtype}
  \UseMicrotypeSet[protrusion]{basicmath} % disable protrusion for tt fonts
}{}
\makeatletter
\@ifundefined{KOMAClassName}{% if non-KOMA class
  \IfFileExists{parskip.sty}{%
    \usepackage{parskip}
  }{% else
    \setlength{\parindent}{0pt}
    \setlength{\parskip}{6pt plus 2pt minus 1pt}}
}{% if KOMA class
  \KOMAoptions{parskip=half}}
\makeatother
\usepackage{graphicx}
\makeatletter
\newsavebox\pandoc@box
\newcommand*\pandocbounded[1]{% scales image to fit in text height/width
  \sbox\pandoc@box{#1}%
  \Gscale@div\@tempa{\textheight}{\dimexpr\ht\pandoc@box+\dp\pandoc@box\relax}%
  \Gscale@div\@tempb{\linewidth}{\wd\pandoc@box}%
  \ifdim\@tempb\p@<\@tempa\p@\let\@tempa\@tempb\fi% select the smaller of both
  \ifdim\@tempa\p@<\p@\scalebox{\@tempa}{\usebox\pandoc@box}%
  \else\usebox{\pandoc@box}%
  \fi%
}
% Set default figure placement to htbp
\def\fps@figure{htbp}
\makeatother
\setlength{\emergencystretch}{3em} % prevent overfull lines
\providecommand{\tightlist}{%
  \setlength{\itemsep}{0pt}\setlength{\parskip}{0pt}}
\usepackage{color}
\usepackage{framed}
\setlength{\fboxsep}{.8em}

\usepackage{tcolorbox}

\tcbset{
  mystyle/.style={
    before skip=-10pt, % Adjust this value
    after skip=-10pt   % Adjust this value
  }
}
\newtcolorbox{mynoteblock}[1][]{mystyle, #1}

\newtcolorbox{blackbox}{
  colback=black,
  colframe=orange,
  coltext=white,
  boxsep=5pt,
  arc=4pt}
  
\newtcolorbox{ltbluebox}{
  colback=blue!5!white,
  colframe=blue!75!black,
  boxsep=5pt,
  arc=4pt}
  
\newtcolorbox{ltgreenbox}{
  colback=green!5!white,
  colframe=blue!75!black,
  boxsep=5pt,
  arc=4pt}
  
\newtcolorbox{ltpurplebox}{
  colback=purple!5!white,
  colframe=blue!75!black,
  boxsep=5pt,
  arc=4pt}
  


\newenvironment{cols}[1][]{}{}

\newenvironment{col}[1]{\begin{minipage}[t]{#1}\ignorespaces\vspace{0pt}}{%
\end{minipage}
\ifhmode\unskip\fi
\aftergroup\useignorespacesandallpars}

\def\useignorespacesandallpars#1\ignorespaces\fi{%
#1\fi\ignorespacesandallpars}

\makeatletter
\def\ignorespacesandallpars{%
  \@ifnextchar\par
    {\expandafter\ignorespacesandallpars\@gobble}%
    {}%
}
\makeatother

\usepackage{awesomebox}
\setlength{\aweboxvskip}{1mm}


\pagenumbering{gobble}

\usepackage{enumitem}
\setlist{itemsep=1pt, parsep=0pt} %# Adjust vertical spacing
\setlist[itemize,1]{leftmargin=*, labelsep=0.3em} %# Adjust horizontal spacing
\setlist[enumerate,1]{leftmargin=*, labelsep=0.3em} %# Adjust horizontal spacing

\usepackage{bookmark}
\IfFileExists{xurl.sty}{\usepackage{xurl}}{} % add URL line breaks if available
\urlstyle{same}
\hypersetup{
  colorlinks=true,
  linkcolor={Maroon},
  filecolor={Maroon},
  citecolor={Blue},
  urlcolor={blue},
  pdfcreator={LaTeX via pandoc}}

\author{}
\date{\vspace{-2.5em}}

\begin{document}

\vspace{-1cm}

\section{Catch/Quota Setting Step 2: Conduct
Assessment}\label{catchquota-setting-step-2-conduct-assessment}

\begin{noteblock}
Catch and Quota Setting is a multistep process involving decisions by
the Northeast Regional Coordinating Committee (NRCC; including the Mid
Atlantic and New England Councils, Atlantic States Marine Fisheries
Commission, and NOAA's Greater Atlantic Regional Fisheries Office and
Northeast Fisheries Science Center), assessment scientists, assessment
peer reviewers, the Council's Scientific and Statistical Committee
(SSC), and the Council. Information for each stock and fishery is
reviewed and synthesized to determine catch and quota levels that meet
management objectives.

\end{noteblock}

\begin{itemize}
\tightlist
\item
  \textbf{Key question: What additional indicators could inform
  decisions and analyses at the Conduct Assessment step in this process?
  Could habitat, ecosystem, or socioeconomic indicators improve the
  process?}
\item
  \textbf{Scope: Should we develop a general process for using
  indicators in the Conduct Assessment step of the catch/quota setting
  process, or focus on a specific application for a single species?
  Which species?}
\end{itemize}

\begin{cols}

\begin{col}{0.49\textwidth}

\begin{ltbluebox}

\subsection{What information is used
now?}\label{what-information-is-used-now}

\begin{ltpurplebox}
\textbf{2. Conduct Assessment}

Assessment scientists assemble and analyze fishery dependent and
independent data and develop models to synthesize the data and produce
outputs meeting the assessment ToRs.

\end{ltpurplebox}

Research Track (RT) Assessment working groups bring in appropriate
expertise to cover the ToRs over a \textasciitilde2 year period.
Ecosystem information is summarized under ToR 1, which may produce an
Ecosystem Socioeconomic Profile (ESP). Assessment leads work with the
ToR 1 team to evaluate whether information can be included in the
assessment model or to inform parameter estimates or structural model
assumptions. \emph{RTs are indefinitely paused so the information
pathway will need to change to introduce new information.}

\vspace{1ex}

ESPs have been produced for 5 stocks during research track assessments
under ToR 1: \href{https://doi.org/10.25923/ez9g-af05}{black sea bass},
\href{https://doi.org/10.25923/s3xc-ta06}{bluefish},
\href{https://doi.org/10.25923/8g2c-0447}{golden tilefish}, longfin
inshore squid (in progress), and
\href{https://doi.org/10.25923/f87x-mj55}{shortfin squid}.

\vspace{1ex}

Management Track (MT) assessments are conducted by the assessment lead
with less formal support in a few month period. MT assessments are used
for management advice while RT assessments are not--therefore ecosystem
information is most likely to be included in the catch quota setting
decision through the assessment if it is approved from the RT assessment
and included in or presented along with the MT assessment to be used at
a later point in the process such as the SSC ABC decision (Catch/Quota
Setting Step 4).

\vspace{2ex}

\textbf{Ecosystem indicators have been integrated into the operational
black sea bass assessment (winter bottom temperature, northern area
influences expected recruitment) and into research track models for
black sea bass and bluefish (forage fish index influences recreational
index catchability). This integration was based on ESPs and accepted
through peer review (Catch/Quota Setting Step 3).}

\end{ltbluebox}

\end{col}

\begin{col}{0.02\textwidth}
~

\end{col}

\begin{col}{0.49\textwidth}

\begin{ltgreenbox}

\subsection{Information sources for
indicators}\label{information-sources-for-indicators}

RT Assessment ToR 1:
\href{https://www.fisheries.noaa.gov/new-england-mid-atlantic/ecosystems/ecosystem-and-socioeconomic-profile-development-and-reports}{Ecosystem
and Socioeconomic Profiles} conceptual models and indicators;
\href{https://www.fisheries.noaa.gov/new-england-mid-atlantic/ecosystems/state-ecosystem-reports-northeast-us-shelf}{State
of the Ecosystem} and
\href{https://www.mafmc.org/s/2_MAB_RiskAssess_2024.pdf}{Risk
Assessment} indicators; The fish and invertebrate
\href{https://journals.plos.org/plosone/article?id=10.1371/journal.pone.0146756}{Climate
Vulnerability Analysis} sensitivity, exposure and vulnerability rankings

\vspace{1ex}

RT ToR 3 or MT ToR 2: Can include ecosystem derived information
(e.g.~seabird diet derived recruitment index) from: *
\href{https://www.fisheries.noaa.gov/inport/item/22549}{NEFSC Surveys} *
\href{https://fwdp.shinyapps.io/tm2020/}{NEFSC Food Habits} *
\href{https://www.vims.edu/research/units/programs/multispecies_fisheries_research/neamap/}{NEAMAP
Survey and Food Habits} *
\href{https://www.fisheries.noaa.gov/new-england-mid-atlantic/commercial-fishing/quota-monitoring-greater-atlantic-region}{GARFO
Catch and Monitoring System} *
\href{https://safis.accsp.org/accsp_prod/f?p=1490:1:6442641492940:::::}{Atlantic
Coastal Cooperative Statistics Programa} *
\href{https://www.fisheries.noaa.gov/recreational-fishing-data/introduction-marine-recreational-information-program-data}{Marine
Recreational Information Program} *
\href{https://data.marine.copernicus.eu/product/GLOBAL_MULTIYEAR_PHY_001_030/description}{GLORYS
ocean reanalysis}

\vspace{1ex}

RT ToR 5 or MT ToR 4: Ecosystem information on changing productivity for
the system or the stock can apply here to determine how long a time
series is the baseline for reference points. MT ToR 4 notes that simple
indicators and metrics of stock status should also be included. Multiple
Essential Fish Habitat sources: *
\href{https://www.fisheries.noaa.gov/resource/map/essential-fish-habitat-mapper}{EFH
Habitat Mapper} * \href{https://nrha.shinyapps.io/dataexplorer/}{NRHA
Data Explorer} *
\href{https://fishmaps.shinyapps.io/FishingEffectsDatabase/}{Fishing
Gear Effects Database}

\vspace{1ex}

RT ToR 6 or MT ToR 5: Ecosystem information on recent productivity for
the system or the stock from sources above can apply here to determine
the basis for short term projections.
\href{https://psl.noaa.gov/cefi_portal/}{MOM6 regional ocean
projections} could also inform short term projections.

\vspace{1ex}

RT ToRs 7 and 8: Identify research to address the ecosystem drivers and
develop ecosystem-based backup methods based on sources above

\begin{ltpurplebox}

\subsection{Information gaps}\label{information-gaps}

Ecosystem and Socioeconomic Profiles evaluate and select stock specific
ecosystem and socioeconomic indicators that can be useful across many
steps of the Catch/Quota setting process; these are currently
unavailable for summer flounder, scup, mackerel, butterfish, surfclam,
ocean quahog, blueline tilefish, spiny dogfish, and monkfish.

\end{ltpurplebox}

\end{ltgreenbox}

\end{col}

\end{cols}

\end{document}
